\section{习题课 \#7（10月25日）}

\subsection{Sylvester 秩不等式及等号条件}

设
\begin{equation}
 A\in\F^{m\times n},
 \qquad
 B\in\F^{n\times p}.
\end{equation}

\begin{theorem}[Sylvester 秩不等式]
有
\begin{equation}
 \rank(AB)\geq\rank A+\rank B-n.
\end{equation}
并且等号成立当且仅当
\begin{equation}
 \Ker A\subseteq\Image B.
\end{equation}
\end{theorem}

\begin{proof}
将 $A$ 限制到 $\Image B$ 上。其像为 $\Image(AB)$，核为 $\Ker A\cap\Image B$，故
\begin{equation}
 \rank(AB)
 =\rank B-\dim(\Ker A\cap\Image B).
\end{equation}
因为
\begin{equation}
 \dim(\Ker A\cap\Image B)\leq\dim\Ker A=n-\rank A,
\end{equation}
即得不等式。等号成立恰好要求
\begin{equation}
 \dim(\Ker A\cap\Image B)=\dim\Ker A,
\end{equation}
也就是 $\Ker A\subseteq\Image B$。
\end{proof}

\begin{corollary}
下列条件等价：
\begin{enumerate}
  \item $\rank(AB)=\rank A+\rank B-n$；
  \item $\Ker A\subseteq\Image B$；
  \item 存在 $X\in\F^{n\times m}$ 与 $Y\in\F^{p\times n}$，使
  \begin{equation}
   XA+BY=I_n.
  \end{equation}
\end{enumerate}
\end{corollary}

\begin{proof}
第三个条件意味着每个 $v\in\Ker A$ 满足 $v=BYv\in\Image B$。反之，若 $\Ker A\subseteq\Image B$，取直和分解
\begin{equation}
 \F^n=U\oplus\Ker A.
\end{equation}
映射 $A|_U$ 单射。可构造 $X$ 使 $XA$ 是到 $U$ 的投影，再构造 $Y$ 使 $BY$ 是到 $\Ker A$ 的投影，于是二者之和为恒等映射。
\end{proof}

\subsection{满秩分解与秩标准形}

\begin{theorem}[满秩分解]
若 $A\in\F^{m\times n}$ 且 $\rank A=r$，则存在
\begin{equation}
 P\in\F^{m\times r},
 \qquad
 Q\in\F^{r\times n}
\end{equation}
使
\begin{equation}
 A=PQ,
 \qquad
 \rank P=\rank Q=r.
\end{equation}
\end{theorem}

\begin{proof}
从 $A$ 中选取 $r$ 个极大线性无关列组成 $P$。$A$ 的每一列都可唯一表示为 $P$ 的列向量的线性组合，把这些坐标作为 $Q$ 的列便有 $A=PQ$。由于 $A$ 的秩为 $r$，$Q$ 必满行秩。
\end{proof}

\begin{proposition}[满秩分解的唯一性]
若
\begin{equation}
 A=P_1Q_1=P_2Q_2
\end{equation}
都是满秩分解，则存在唯一 $S\in\GL_r(\F)$ 使
\begin{equation}
 P_2=P_1S,
 \qquad
 Q_2=S^{-1}Q_1.
\end{equation}
\end{proposition}

\begin{proof}
$P_1,P_2$ 的列空间都等于 $\Image A$，故 $P_2=P_1S$。两组列都是基，所以 $S$ 可逆。代回并利用 $P_1$ 满列秩即可得到 $Q_2=S^{-1}Q_1$。
\end{proof}

\begin{corollary}[秩标准形]
存在可逆矩阵 $U\in\GL_m(\F)$、$V\in\GL_n(\F)$，使
\begin{equation}
 UAV=\begin{pmatrix}
 I_r&0\\
 0&0
 \end{pmatrix}.
\end{equation}
因此两个同型矩阵等价当且仅当它们的秩相同。
\end{corollary}

\begin{remark}
若经过行列置换后，$A$ 的左上角 $r\times r$ 子矩阵 $A_{11}$ 可逆，写成
\begin{equation}
 A=\begin{pmatrix}
 A_{11}&A_{12}\\
 A_{21}&A_{22}
 \end{pmatrix},
\end{equation}
则 $\rank A=r$ 当且仅当
\begin{equation}
 A_{22}=A_{21}A_{11}^{-1}A_{12}.
\end{equation}
这就是 Schur 补为零的情形。
\end{remark}

\subsection{矩阵和的秩}

\begin{theorem}
对同型矩阵 $A,B$，
\begin{equation}
 \rank(A+B)\leq\rank A+\rank B.
\end{equation}
等号成立当且仅当
\begin{equation}
 \Image A\cap\Image B=\{0\}
\end{equation}
且
\begin{equation}
 \Image A^T\cap\Image B^T=\{0\}.
\end{equation}
复数情形可将转置理解为行空间，而不必使用共轭转置。
\end{theorem}

\begin{proof}
因为
\begin{equation}
 \Image(A+B)\subseteq\Image A+\Image B,
\end{equation}
先得上界。若等号成立，则
\begin{equation}
 \rank(A+B)=\dim(\Image A+\Image B)
\end{equation}
迫使列空间之交为零；对转置应用同一论证得到行空间之交为零。

反之，若两种交都为零，可选取适当的可逆行、列变换，使
\begin{equation}
 A\sim
 \begin{pmatrix}
 I_r&0&0\\
 0&0&0\\
 0&0&0
 \end{pmatrix},
 \qquad
 B\sim
 \begin{pmatrix}
 0&0&0\\
 0&I_s&0\\
 0&0&0
 \end{pmatrix},
\end{equation}
且两者使用相同的外侧变换。于是 $A+B$ 的秩为 $r+s$。
\end{proof}

\begin{corollary}
对同型矩阵 $A,B$，以下三条中的任意两条推出第三条：
\begin{equation}
 \rank(A+B)=\rank A+\rank B,
\end{equation}
\begin{equation}
 \rank(A^T,B^T)=\rank A+\rank B,
\end{equation}
以及相应列空间、行空间直和条件。
\end{corollary}

\subsection{Binet--Cauchy 公式与伴随矩阵}

\begin{theorem}[Binet--Cauchy 公式]
设 $A\in\F^{m\times n}$、$B\in\F^{n\times m}$，且 $m\leq n$。对 $m$ 元指标集
\begin{equation}
 I=\{i_1<\cdots<i_m\}\subseteq\{1,\ldots,n\},
\end{equation}
记 $A_{:,I}$ 为取 $A$ 的这些列所得方阵，$B_{I,:}$ 为取 $B$ 的这些行所得方阵。则
\begin{equation}
 \det(AB)=\sum_{|I|=m}\det(A_{:,I})\det(B_{I,:}).
\end{equation}
\end{theorem}

\begin{proof}
把 $AB$ 的每一列写成 $A$ 的列向量的线性组合，并对行列式使用多重线性性。若所选指标有重复，相应行列式为零；将互异指标按递增顺序重排后，两个置换符号相消，便得到上述求和式。
\end{proof}

\begin{corollary}[Gram 行列式]
若 $A\in\R^{n\times m}$ 且 $n\geq m$，则
\begin{equation}
 \det(A^TA)=\sum_{|I|=m}\det(A_{I,:})^2\geq0.
\end{equation}
并且 $\det(A^TA)>0$ 当且仅当 $A$ 满列秩。
\end{corollary}

\begin{proposition}
对任意同阶方阵 $A,B$，有
\begin{equation}
 \adj(AB)=\adj(B)\adj(A).
\end{equation}
\end{proposition}

\begin{proof}
当 $A,B$ 都可逆时，
\begin{equation}
 \adj(AB)=\det(AB)(AB)^{-1}
 =\det B\,B^{-1}\det A\,A^{-1}
 =\adj(B)\adj(A).
\end{equation}
两边的每个矩阵元都是 $A,B$ 各元素的多项式。把 $A,B$ 分别替换为 $A+tI,B+sI$，在除有限多个参数值外二者可逆，于是多项式恒等式对所有 $A,B$ 成立。
\end{proof}

\subsection{低秩修正与 Schur 补}

\begin{theorem}[对角矩阵加秩一矩阵]
设
\begin{equation}
 D=\diag(a_1,\ldots,a_n)
\end{equation}
可逆，且
\begin{equation}
 1+\ones^TD^{-1}\ones\neq0.
\end{equation}
则
\begin{equation}
 (D+\ones\ones^T)^{-1}
 =D^{-1}-
 \frac{D^{-1}\ones\ones^TD^{-1}}
 {1+\ones^TD^{-1}\ones}.
\end{equation}
\end{theorem}

\begin{proof}
将右端记为 $R$，直接计算
\begin{equation}
 (D+\ones\ones^T)R=I_n.
\end{equation}
这也是 Sherman--Morrison 公式的特例。
\end{proof}

\begin{theorem}[Schur 补与分块逆]
设
\begin{equation}
 M=\begin{pmatrix}A&B\\ C&D\end{pmatrix},
\end{equation}
其中 $A$ 可逆，并记
\begin{equation}
 S=D-CA^{-1}B.
\end{equation}
则
\begin{equation}
 \begin{pmatrix}I&0\\-CA^{-1}&I\end{pmatrix}
 M
 \begin{pmatrix}I&-A^{-1}B\\0&I\end{pmatrix}
 =\begin{pmatrix}A&0\\0&S\end{pmatrix}.
\end{equation}
因此 $M$ 可逆当且仅当 $S$ 可逆，并且
\begin{equation}
 M^{-1}=
 \begin{pmatrix}
 A^{-1}+A^{-1}BS^{-1}CA^{-1}&-A^{-1}BS^{-1}\\
 -S^{-1}CA^{-1}&S^{-1}
 \end{pmatrix}.
\end{equation}
\end{theorem}

\begin{corollary}
若 $M=M^T$ 且 $A$ 正定，则
\begin{equation}
 M\text{ 正定}
 \quad\Longleftrightarrow\quad
 S=D-B^TA^{-1}B\text{ 正定}.
\end{equation}
\end{corollary}

\subsection{Vandermonde 矩阵的逆}

令 $x_1,\ldots,x_n$ 两两不同，并取标准 Vandermonde 矩阵
\begin{equation}
 V=(x_i^{j-1})_{i,j=1}^n.
\end{equation}
定义 Lagrange 基函数
\begin{equation}
 \ell_i(t)=\prod_{k\neq i}\frac{t-x_k}{x_i-x_k}
 =\sum_{j=1}^n c_{ji}t^{j-1}.
\end{equation}
由于 $\ell_i(x_k)=\delta_{ik}$，有
\begin{equation}
 V(c_{ji})_{j,i=1}^n=I_n.
\end{equation}

\begin{theorem}
Vandermonde 矩阵的逆满足
\begin{equation}
 (V^{-1})_{ji}=[t^{j-1}]\ell_i(t),
\end{equation}
即第 $i$ 列由第 $i$ 个 Lagrange 基多项式的系数组成。显式地，
\begin{equation}
 (V^{-1})_{ji}
 =\frac{(-1)^{n-j}
 e_{n-j}(x_1,\ldots,\widehat{x_i},\ldots,x_n)}
 {\prod_{k\neq i}(x_i-x_k)},
\end{equation}
其中 $e_k$ 是初等对称多项式，帽号表示删去该变量。
\end{theorem}

\begin{remark}[单位根情形]
取
\begin{equation}
 \omega=\e^{-2\pi\ii/n},
 \qquad
 F=(\omega^{jk})_{j,k=0}^{n-1},
\end{equation}
则几何级数给出
\begin{equation}
 F^*F=nI_n,
 \qquad
 F^{-1}=\frac1nF^*.
\end{equation}
这正是 Vandermonde 逆在单位根节点上的特殊简化。
\end{remark}

\subsection{行列式交换恒等式与矩阵逆公式}

\begin{theorem}[Sylvester 行列式恒等式]
若 $A\in\F^{m\times n}$、$B\in\F^{n\times m}$，则
\begin{equation}
 \det(I_m+AB)=\det(I_n+BA).
\end{equation}
更一般地，作为关于 $\lambda$ 的多项式，
\begin{equation}
 \det(\lambda I_m-AB)
 =\lambda^{m-n}\det(\lambda I_n-BA)
\end{equation}
在 $m\geq n$ 时成立；若 $m<n$，交换 $m,n$ 后理解该式。
\end{theorem}

\begin{proof}
对分块矩阵
\begin{equation}
 \begin{pmatrix}I_m&A\\-B&I_n\end{pmatrix}
\end{equation}
分别以左上角和右下角作 Schur 补，即得第一式。第二式先在 $\lambda\neq0$ 时提出 $\lambda$ 并应用第一式，再由多项式恒等性延拓到所有 $\lambda$。
\end{proof}

\begin{corollary}
$AB$ 与 $BA$ 的非零特征值连同代数重数完全相同。
\end{corollary}

\begin{proposition}[Neumann 级数]
若某个次乘法矩阵范数满足 $\norm{A}<1$，则
\begin{equation}
 (I-A)^{-1}=I+A+A^2+\cdots,
\end{equation}
级数在该范数下绝对收敛。更一般地，只要谱半径 $\rho(A)<1$，该结论仍成立。
\end{proposition}

\begin{proposition}[推入恒等式]
若相关逆存在，则
\begin{equation}
 (I-AB)^{-1}=I+A(I-BA)^{-1}B.
\end{equation}
等价地，
\begin{equation}
 (I-AB)^{-1}A=A(I-BA)^{-1}.
\end{equation}
\end{proposition}

\begin{theorem}[Woodbury 公式]
若各逆均存在，则
\begin{equation}
 (A+UCV)^{-1}
 =A^{-1}-A^{-1}U
 (C^{-1}+VA^{-1}U)^{-1}
 VA^{-1}.
\end{equation}
当 $C=1$ 且 $U=u,V=v^T$ 时，得到 Sherman--Morrison 公式
\begin{equation}
 (A+uv^T)^{-1}
 =A^{-1}-\frac{A^{-1}uv^TA^{-1}}
 {1+v^TA^{-1}u}.
\end{equation}
\end{theorem}

\subsection{DFT、循环矩阵与快速算法}

取 $\omega=\e^{-2\pi\ii/n}$，离散 Fourier 矩阵为
\begin{equation}
 F=(\omega^{jk})_{j,k=0}^{n-1}.
\end{equation}
对向量 $x\in\C^n$，其离散 Fourier 变换和逆变换分别为
\begin{equation}
 \widehat{x}=Fx,
 \qquad
 x=\frac1nF^*\widehat{x}.
\end{equation}

\begin{theorem}[循环矩阵的对角化]
设 $C$ 的第一列为 $c=(c_0,\ldots,c_{n-1})^T$，其余列由循环移位得到。则
\begin{equation}
 C=F^{-1}\diag(Fc)F.
\end{equation}
因此循环卷积满足卷积定理
\begin{equation}
 F(x\ast y)=(Fx)\odot(Fy),
\end{equation}
其中 $\odot$ 表示逐分量乘积。
\end{theorem}

\begin{proof}
Fourier 向量
\begin{equation}
 v_k=(1,\omega^k,\omega^{2k},\ldots,\omega^{(n-1)k})^T
\end{equation}
是每个循环移位算子的特征向量，因而也是任意循环矩阵的特征向量。相应特征值正是 $Fc$ 的各分量。
\end{proof}

\begin{remark}
直接计算 DFT 需要 $O(n^2)$ 次运算。快速 Fourier 变换利用 $n$ 的因子分解递归拆分偶、奇指标，将复杂度降为 $O(n\log n)$。把单位根换成有限域中的适当本原根便得到数论变换，可用于避免浮点误差的多项式乘法。
\end{remark}
